\documentclass[11pt]{amsart}
\usepackage[utf8]{inputenc}
\usepackage{amsmath,amssymb,amsthm,mathtools}
\usepackage[margin=1.1in]{geometry}
\usepackage{hyperref}

\newtheorem{theorem}{Theorem}
\newtheorem{proposition}[theorem]{Proposition}
\newtheorem{lemma}[theorem]{Lemma}
\newtheorem{corollary}[theorem]{Corollary}
\newtheorem{conjecture}[theorem]{Conjecture}
\theoremstyle{definition}
\newtheorem{definition}[theorem]{Definition}
\newtheorem{remark}[theorem]{Remark}
\newtheorem{question}[theorem]{Question}

\newcommand{\R}{\mathbb{R}}
\newcommand{\conv}{\operatorname{conv}}
\newcommand{\vol}{\operatorname{vol}}
\newcommand{\C}{\mathbb{C}}
\newcommand{\re}{\operatorname{Re}}

\title[Signature volume formula and cyclic curves]{On the Am\'endola--Lee--Meroni conjecture:\\
the signature volume formula on all subpaths does not characterize cyclic curves}
\author{Paata Ivanisvili}
\date{October 26, 2026}

\begin{document}
\begin{abstract}
Am\'endola, Lee and Meroni proved that the volume of the convex hull of a cyclic curve
$x:[0,1]\to\R^d$ equals its top-level alternating path signature $\alpha^{(d)}(x)$, and they conjectured
(GSI 2023, Conjecture 3.11) that cyclic curves are exactly the curves for which this identity holds for every subpath.
The literal statement fails trivially for non-cyclic curves inside a hyperplane. We show that full-dimensionality does not save it: the conjecture is false in every
dimension $d\ge 2$, even for full-dimensional piecewise linear curves. In the plane we prove that the conjecture is true for injective curves (simple arcs and simple
closed curves). In every dimension $d\ge3$ it fails even for injective full-dimensional curves: a cone over a logarithmic spiral started at its pole,
with a specific growth rate, satisfies the identity on every subpath but is not cyclic. We also record, with a Lipschitz bound, the known fact that a
full-dimensional curve whose ordered $(d+1)$-point determinants are all nonnegative is cyclic.
\end{abstract}
% --- AI disclosure (site) ---
\noindent\begin{center}\small\itshape
This proof was generated on October 27, 2026, by Grok Bot using Cursor’s agents, guided by several prompts from Paata Ivanisvili
\end{center}\medskip
% --- end AI disclosure ---
\maketitle

\section{Setting and statement}

All paths are Lipschitz maps $x:[a,b]\to\R^d$; $x'$ exists a.e. For $k\ge1$ let
\[
S^{(k)}(x)=\int_{a<t_1<\dots<t_k<b} x'(t_1)\wedge\dots\wedge x'(t_k)\,dt\in\Lambda^k\R^d,\qquad S^{(0)}(x)=1,
\]
and $\alpha^{(d)}(x)=\frac1{d!}\langle S^{(d)}(x),e_1\wedge\dots\wedge e_d\rangle=\frac1{d!}\int_{a<t_1<\dots<t_d<b}\det(x'(t_1),\dots,x'(t_d))\,dt$.
This is the normalization of \cite{ALM}. For $d=2$,
\[
\alpha^{(2)}(x|_{[a,b]})=A(a,b):=\tfrac12\int_a^b\det\big(x(t)-x(a),x'(t)\big)\,dt,
\]
which is the L\'evy area: the signed area enclosed by $x|_{[a,b]}$ followed by the chord from $x(b)$ to $x(a)$.
For points $p_0,\dots,p_d\in\R^d$ put $\Delta(p_0,\dots,p_d)=\det\begin{pmatrix}1&\cdots&1\\ p_0&\cdots&p_d\end{pmatrix}$.

\begin{definition}[\cite{ALM}]
A path $x:[0,1]\to\R^d$ is \emph{$d$-order} if $\Delta(x(t_0),\dots,x(t_d))>0$ for all $0\le t_0<\dots<t_d\le1$.
It is \emph{cyclic} if it is a uniform limit of $d$-order paths. (\cite{ALM}, Def.~2.6, allow limits in the Lipschitz or, equivalently they say, the
uniform topology. We use uniform limits. The counterexamples below are not weakly ordered, so they are not limits in either sense. The positive results produce
uniform limits with bounded Lipschitz constants.)
We call $x$ \emph{weakly $d$-ordered} if $\Delta(x(t_0),\dots,x(t_d))\ge0$ for all $t_0<\dots<t_d$.
\end{definition}

Weak order is preserved under pointwise limits, so cyclic paths are weakly $d$-ordered.

\begin{theorem}[\cite{ALM}, Thms.\ 3.4--3.5]\label{thm:ALM}
If $x$ is $d$-order, or a uniform limit of $d$-order paths with uniformly bounded Lipschitz constants, then
$\vol_d(\conv x([0,1]))=\alpha^{(d)}(x)$.
\end{theorem}

We say that $x$ satisfies \textbf{(H)} if $\vol_d(\conv x([a,b]))=\alpha^{(d)}(x|_{[a,b]})$ for all $0\le a<b\le1$.
Cyclic Lipschitz paths satisfy (H): they are weakly $d$-ordered, so Proposition~\ref{prop:smooth} below applies (and if $x([a,b])$ lies in a hyperplane, both sides vanish).

\begin{conjecture}[\cite{ALM}, Conjecture 3.11]\label{conj:ALM}
A Lipschitz path $x:[0,1]\to\R^d$ is cyclic if and only if it satisfies (H).
\end{conjecture}

\begin{theorem}\label{thm:counter}
For every $d\ge2$ there is a piecewise linear path $x:[0,1]\to\R^d$ that satisfies (H), whose image is not contained in any
affine hyperplane, and that is not weakly $d$-ordered; in particular it is not cyclic. Hence Conjecture~\ref{conj:ALM} is false for every $d\ge2$.
\end{theorem}

\begin{theorem}\label{thm:inj}
For every $d\ge3$ there is an injective Lipschitz path $x:[0,1]\to\R^d$, not contained in any affine hyperplane, that satisfies (H) and is not weakly $d$-ordered.
In particular Conjecture~\ref{conj:ALM} fails for injective full-dimensional paths in every dimension $d\ge3$.
\end{theorem}

\begin{theorem}\label{thm:planar}
Let $x:[0,1]\to\R^2$ be Lipschitz, injective on $[0,1)$ and on $(0,1]$, so that $x$ is a simple arc or a simple closed curve.
The following are equivalent: (a) $x$ satisfies (H); (b) $x$ is weakly $2$-ordered; (c) $x$ is cyclic.
In particular Conjecture~\ref{conj:ALM} holds for such curves.
\end{theorem}

\begin{proposition}\label{prop:smooth}
Let $x:[0,1]\to\R^d$ be continuous and weakly $d$-ordered, with $x([0,1])$ not contained in an affine hyperplane.
Then there are $C^\infty$ $d$-order paths $x_\varepsilon\to x$ uniformly as $\varepsilon\to0$.
If $x$ is Lipschitz, they can be chosen with $\sup_\varepsilon\operatorname{Lip}(x_\varepsilon)<\infty$, and then $x$ satisfies (H).
In particular $x$ is cyclic.
\end{proposition}

So a full-dimensional path is cyclic if and only if it is weakly $d$-ordered. \cite{ALM} state this as a remark. In substance it is the classical fact that weak
Tchebycheff systems are limits of Tchebycheff systems, obtained by smoothing with the Gauss kernel (\cite{K}, Ch.~3; \cite{Z}). Here $(1,x_1,\dots,x_d)$ is a weak Tchebycheff
system exactly when $x$ is weakly $d$-ordered. The new point in Proposition~\ref{prop:smooth} is only the Lipschitz bound. The counterexamples of Theorems~\ref{thm:counter} and
\ref{thm:inj} satisfy (H) but are not weakly $d$-ordered. Theorem~\ref{thm:planar} shows that in the plane every counterexample must revisit points, and
Theorem~\ref{thm:inj} shows that this is not so for $d\ge3$.

\begin{question}\label{q:inj}
Are there injective counterexamples in $d\ge3$ that are piecewise linear or smooth? The examples of Theorem~\ref{thm:inj} wind infinitely often around one point.
What does (H) characterize?
\end{question}

\section{Two algebraic tools}

\begin{lemma}[Chen's identity]\label{lem:chen}
If $x$ is the concatenation of $y$ (on $[a,b]$) and $z$ (on $[b,c]$), then
$S^{(k)}(x)=\sum_{i+j=k}S^{(i)}(y)\wedge S^{(j)}(z)$.
\end{lemma}
\begin{proof}
Split $\{a<t_1<\dots<t_k<c\}$ according to the number $i$ of the $t_\nu$ lying in $[a,b]$. The piece with index $i$ is
$\{a<t_1<\dots<t_i<b\}\times\{b<t_{i+1}<\dots<t_k<c\}$ up to a null set. On it the integrand factors, by multilinearity of $\wedge$.
\end{proof}

For a segment $y$, $S^{(j)}(y)=0$ for $j\ge2$ because $v\wedge v=0$.
For a path contained in an affine hyperplane $p+W$, $S^{(j)}(y)\in\Lambda^jW$, so $S^{(d)}(y)=0$; both sides of (H) then vanish.

\begin{lemma}[Cone lemma]\label{lem:cone}
Let $Y:[0,c]\to\R^{d-1}\times\{0\}$ be Lipschitz with $Y(0)=P$. Let $\epsilon=(-1)^d$ and $h>0$. Let $X$ be the segment from
$P+\epsilon h e_d$ to $P$, and let $Z=X*Y$. Then
\[
\vol_d(\conv Z)=\tfrac hd\vol_{d-1}(\conv Y),\qquad \alpha^{(d)}(Z)=\tfrac hd\,\alpha^{(d-1)}(Y),
\]
where the right-hand sides are computed in $\R^{d-1}$. Moreover, for $t_0$ in the interior of the parameter interval of $X$,
with $Z(t_0)=P+\epsilon h'e_d$, and $t_1<\dots<t_d$ in the parameter interval of $Y$,
$\Delta(Z(t_0),Z(t_1),\dots,Z(t_d))=h'\,\Delta_{\R^{d-1}}(Y(t_1),\dots,Y(t_d))$.
\end{lemma}
\begin{proof}
Since $P\in\conv Y$, $\conv Z=\conv(\{P+\epsilon he_d\}\cup\conv Y)$, a pyramid of height $h$ over the $(d-1)$-dimensional base $\conv Y$.
By Lemma~\ref{lem:chen} and the remarks after it, $S^{(d)}(Z)=S^{(1)}(X)\wedge S^{(d-1)}(Y)=-\epsilon h\,e_d\wedge S^{(d-1)}(Y)$.
Moving $e_d$ to the last position costs $(-1)^{d-1}$, and $(-1)^{d-1}(-\epsilon)=1$. Hence
$\langle S^{(d)}(Z),e_1\wedge\dots\wedge e_d\rangle=h\langle S^{(d-1)}(Y),e_1\wedge\dots\wedge e_{d-1}\rangle$; divide by $d!=d\cdot(d-1)!$.
For the determinant, write $Z(t_i)-Z(t_0)=(w_i,-\epsilon h')$ with $w_i=Y(t_i)-P\in\R^{d-1}$. Then
$\det(Z(t_1)-Z(t_0),\dots,Z(t_d)-Z(t_0))=(-\epsilon h')(-1)^{d-1}\Delta_{\R^{d-1}}(w_1,\dots,w_d)$.
The sign is $(-1)^{d+1}(-1)^{d-1}=1$, and $\Delta$ is translation invariant.
\end{proof}

\section{Counterexamples: proof of Theorem~\ref{thm:counter}}

\subsection*{The planar example}
Let $\Gamma:\R\to\R^2$ be the $4$-periodic, unit-speed, counterclockwise parametrization of $\partial Q$, $Q=[0,1]^2$, with
$\Gamma(0)=(0,0)$: $\Gamma(t)=(t,0)$ on $[0,1]$, $(1,t-1)$ on $[1,2]$, $(3-t,1)$ on $[2,3]$, $(0,4-t)$ on $[3,4]$. Define
\[
x=\Gamma|_{[0,9/2]}.
\]
This is one counterclockwise lap of the square followed by the segment from $(0,0)$ to $(1/2,0)$ along the first edge.
After rescaling the parameter to $[0,1]$, which changes neither side of (H), we check (H) on each $[a,b]\subset[0,9/2]$.

\emph{Case $b-a\le4$.} Then $x|_{[a,b]}$ is a counterclockwise boundary arc $\sigma$ of the convex polygon $Q$ of length at most
the perimeter. The closed curve $\sigma$ followed by the chord $[x(b),x(a)]$ is the positively oriented boundary of the convex set
$\conv\sigma$, which is the part of $Q$ cut off by the chord on the side of $\sigma$, or $Q$ itself when $b-a=4$.
Hence $A(a,b)=|\conv\sigma|=V(a,b)$, where $V(a,b):=|\conv x([a,b])|$.

\emph{Case $b-a>4$.} Then $0\le a<b-4\le1/2$, so $x([a,b])=\partial Q$ and $V(a,b)=1$. By Lemma~\ref{lem:chen} with the split point $a+4$,
\[
A(a,b)=A(a,a+4)+A(a+4,b)+\tfrac12\det\big(x(a+4)-x(a),x(b)-x(a+4)\big)=1+0+0,
\]
because $x(a+4)=x(a)$ and $x|_{[a+4,b]}$ is a segment.

So (H) holds. However $\Delta(x(0),x(2),x(4.3))=\det((1,1),(0.3,0))=-0.3<0$, so $x$ is not weakly $2$-ordered and hence not cyclic.
Concretely, the line $\{x_1=1/4\}$ is crossed transversally at $t=1/4$, $11/4$ and $17/4$.

\subsection*{Induction on $d$}
Assume $Y_{d-1}:[0,c]\to\R^{d-1}$ is piecewise linear, full-dimensional, satisfies (H), and has an ordered $d$-tuple of parameters
with $\Delta_{\R^{d-1}}<0$. Let $P=Y_{d-1}(0)$ and let $Z=Y_d$ be as in Lemma~\ref{lem:cone} with $h=1$.
Full dimensionality is clear. For a subinterval of $Z$ there are three cases.
If it lies in $X$, both sides of (H) are $0$ (a segment).
If it lies in the domain of $Y_{d-1}$, both sides are $0$ (contained in a hyperplane).
Otherwise it starts at a parameter $t_a$ in the domain of $X$, with $Z(t_a)=P+\epsilon h_ae_d$, $0<h_a\le h$. Then $Z$ restricted to it is a segment from
$P+\epsilon h_ae_d$ to $P$ followed by $Y_{d-1}|_{[0,c']}$, and Lemma~\ref{lem:cone} (with $h_a$ in place of $h$)
together with (H) for $Y_{d-1}$ on $[0,c']$ gives (H) for $Z$.
Finally Lemma~\ref{lem:cone} turns the negative $d$-tuple of $Y_{d-1}$ into a negative $(d+1)$-tuple of $Z$.
This proves Theorem~\ref{thm:counter}. \qed

\begin{remark}
(1) There are also degenerate counterexamples: any non-cyclic path inside an affine hyperplane, for example
$t\mapsto(\sin4\pi t,0)$ in $\R^2$, satisfies (H) trivially.
(2) The planar example can be made $C^{1,1}$. Replace $Q$ by a stadium, the convex hull of two equal discs, and let the tail run along a
straight part of its boundary. The same two cases apply.
(3) In these examples the path revisits points it has already passed. Theorem~\ref{thm:planar} shows that in the plane this is unavoidable. For $d\ge3$ it is not; see Theorem~\ref{thm:inj}.
\end{remark}

\section{Injective counterexamples in $d\ge3$: proof of Theorem~\ref{thm:inj}}

The induction step in \S3 uses (H) for the base $Y_{d-1}$ only on \emph{initial} intervals $[0,c']$: subintervals inside the base lie in a hyperplane, where both sides
of (H) vanish. Theorem~\ref{thm:planar} uses (H) on all subintervals. This leaves room for an injective planar base.

\subsection*{The spiral}
For $k>0$ let $Y(\varphi)=e^{(k+i)\varphi}$ ($\varphi\le0$) in $\C=\R^2$, with $Y(-\infty)=0$. Its length $\int_{-\infty}^0|Y'|=\sqrt{1+k^2}/k$ is finite, so the arc-length
parametrization $y:[0,\ell]\to\R^2$ starting at the pole $y(0)=0$ is an injective Lipschitz path. Let $V(\varphi_b)$ be the area of $\conv Y((-\infty,\varphi_b])$ and $A(\varphi_b)$ the
L\'evy area of $y$ up to $Y(\varphi_b)$. Since $\det(Y,Y')=|Y|^2=e^{2k\varphi}$,
\[
A(\varphi_b)=\tfrac12\int_{-\infty}^{\varphi_b}e^{2k\varphi}\,d\varphi=\frac{e^{2k\varphi_b}}{4k}.
\]
The similarity $w\mapsto e^{(k+i)\varphi_b}w$ maps $Y((-\infty,0])$ onto $Y((-\infty,\varphi_b])$, so $V(\varphi_b)=e^{2k\varphi_b}V(0)$. Hence:

\emph{$y$ satisfies (H) on all initial intervals if and only if $V(0)=\frac1{4k}$.}

\begin{lemma}\label{lem:spiral}
Let $\gamma=\arctan(1/k)\in(0,\frac\pi2)$. The boundary of $\conv Y((-\infty,0])$ consists of the arc $Y([\varphi_t,0])$ and the segment from $Y(0)$ to $Y(\varphi_t)$, which is
tangent to the spiral at $Y(\varphi_t)$. Here $\varphi_t\in(-2\pi,-\pi)$ is the unique solution in $(-2\pi,-\pi-\gamma)$ of $k\sin\varphi+\cos\varphi=e^{k\varphi}$. Consequently
\[
V(0)=\frac{1-e^{2k\varphi_t}}{4k}+\tfrac12e^{k\varphi_t}\sin\varphi_t .
\]
\end{lemma}
\begin{proof}
For a direction $\theta$, the support function is $h(\theta)=\sup_{\varphi\le0}g_\theta(\varphi)$ with $g_\theta(\varphi)=\re(e^{-i\theta}Y(\varphi))=e^{k\varphi}\cos(\varphi-\theta)$.
Since $g_\theta'=e^{k\varphi}\sqrt{1+k^2}\cos(\varphi-\theta+\gamma)$, the local maxima of $g_\theta$ are at $\varphi_n=\theta+\frac\pi2-\gamma+2\pi n$, with values
$e^{k\varphi_n}\sin\gamma$, which increase with $n$. Put $\theta_0=\gamma-\frac\pi2$ and write $\theta=\theta_0+\psi$ with $\psi\in[0,2\pi)$. The largest local maximum in
$(-\infty,0]$ is then at $\varphi_{-1}=\psi-2\pi$, and $h(\theta)=\max\big(g_\theta(0),g_\theta(\psi-2\pi)\big)$. The difference is
$D(\psi)=g_\theta(0)-g_\theta(\psi-2\pi)=\sin(\psi+\gamma)-\sin\gamma\,e^{k(\psi-2\pi)}$.
We have $D(0)=\sin\gamma\,(1-e^{-2\pi k})>0$. On $[\pi-\gamma,2\pi-\gamma]$, $D<0$, since the first term is $\le0$. On $(2\pi-\gamma,2\pi)$,
$D'=\cos(\psi+\gamma)-k\sin\gamma\,e^{k(\psi-2\pi)}>\cos\gamma-k\sin\gamma=0$ and $D(2\pi)=0$, so $D<0$. On $(0,\pi-\gamma)$, $D''<0$. So $D$ has exactly one zero $\psi_t$ in
$(0,2\pi)$, and $\psi_t\in(0,\pi-\gamma)$.
For $\psi\in(\psi_t,2\pi)$ the supporting point is $Y(\psi-2\pi)$, and for $\psi\in[0,\psi_t)$ it is $Y(0)$. So the boundary is the arc $Y([\varphi_t,0])$ with $\varphi_t=\psi_t-2\pi$,
together with the segment from $Y(0)$ to $Y(\varphi_t)$. These two points lie on the same support line, which is perpendicular to the direction $\theta_0+\psi_t$, and at the local
maximum $Y(\varphi_t)$ this line is the tangent line. Expanding $D(\psi_t)=0$ with $\cot\gamma=k$ gives $k\sin\varphi_t+\cos\varphi_t=e^{k\varphi_t}$, and uniqueness follows from
uniqueness of $\psi_t$. The boundary is traversed counterclockwise by the arc followed by the segment. The area formula is $\frac12\oint\det(w,dw)$: the arc contributes
$\frac12\int_{\varphi_t}^0e^{2k\varphi}d\varphi$ and the segment contributes $\frac12\det(Y(0),Y(\varphi_t))=\frac12e^{k\varphi_t}\sin\varphi_t$.
\end{proof}

\begin{lemma}\label{lem:kstar}
Let $\Phi(k)=\frac{2k}{\sqrt{1+k^2}}-e^{k(\arctan(1/k)-2\pi)}$. Then $\Phi(0.19)<0<\Phi(0.2)$, so $\Phi$ has a zero $k^*\in(0.19,0.2)$; numerically $k^*=0.195378813151504\ldots$
For $k=k^*$ we have $\varphi_t=\gamma-2\pi$ and $V(0)=\frac1{4k^*}$.
\end{lemma}
\begin{proof}
$\Phi(0.19)=0.37332\ldots-0.39417\ldots<0$ and $\Phi(0.2)=0.39223\ldots-0.37460\ldots>0$. At $k=k^*$, take $\psi=\gamma$, so $\varphi=\gamma-2\pi$. Using
$e^{k(\gamma-2\pi)}=2k\sin\gamma$ (that is, $\Phi(k^*)=0$, since $\sin\gamma=1/\sqrt{1+k^2}$), we get
$D(\gamma)=\sin2\gamma-2k\sin^2\gamma=2\sin\gamma(\cos\gamma-k\sin\gamma)=0$. Since $\gamma\in(0,\pi-\gamma)$, uniqueness gives $\psi_t=\gamma$, i.e.\ $\varphi_t=\gamma-2\pi$.
Then $\sin\varphi_t=\sin\gamma$ and $e^{k\varphi_t}=2k\sin\gamma$, and Lemma~\ref{lem:spiral} gives
\[
V(0)=\frac1{4k}-\frac{e^{2k\varphi_t}}{4k}+\frac12e^{k\varphi_t}\sin\varphi_t=\frac1{4k}+\frac{e^{k\varphi_t}}{4k}\big(2k\sin\gamma-e^{k\varphi_t}\big)=\frac1{4k}.
\]
\end{proof}

\begin{proof}[Proof of Theorem~\ref{thm:inj}]
Let $k=k^*$ and let $y$ be the arc-length parametrized spiral from the pole, rescaled to $[0,1]$. By Lemma~\ref{lem:kstar} and self-similarity, $y$ satisfies (H) on every
initial interval. Moreover $y$ is not weakly $2$-ordered: for $\varphi_1<\varphi_2$ with $\varphi_2-\varphi_1\in(\pi,2\pi)$,
$\Delta(0,Y(\varphi_1),Y(\varphi_2))=\det(Y(\varphi_1),Y(\varphi_2))=e^{k(\varphi_1+\varphi_2)}\sin(\varphi_2-\varphi_1)<0$.
Now apply Lemma~\ref{lem:cone} with $d=3$, $P=0$ and $h=1$ to obtain $Y_3=X*y$. The three-case argument of the induction step in \S3 uses (H) for the base only on
initial intervals, so it shows that $Y_3$ satisfies (H). The segment meets the plane $\{x_3=0\}$ only at the pole, so $Y_3$ is injective. It is full-dimensional, and by
Lemma~\ref{lem:cone} the negative triple of $y$ (with $t_1=0$, the pole) becomes a negative $4$-tuple of $Y_3$. Iterating the induction step of \S3 gives injective examples in every
$d\ge3$; injectivity is preserved because each new segment meets the previous hyperplane only at its endpoint.
\end{proof}

\noindent\emph{Numerical confirmation} (\texttt{numerics/alm\_spiral\_cone.py}, \texttt{numerics/alm\_spiral\_cone.txt}).
\begin{itemize}
\item The convex hull area of a dense sample of the spiral equals $1/(4k^*)$ to $4\cdot10^{-9}$.
\item For $Y_3$, the hull volume equals $|\alpha^{(3)}|$ to relative accuracy $10^{-7}$ on $33$ random subintervals that start in the segment. The planar subintervals give $0=0$.
\item The ordered $4$-tuple (apex, pole, $Y(-5)$, $Y(-\frac12)$) has $\Delta=-0.334$.
\end{itemize}
The planar spiral alone is not a counterexample in $\R^2$: it violates (H) on some non-initial interval, consistently with Theorem~\ref{thm:planar}.

\section{Proof of Proposition~\ref{prop:smooth}}

For $\varepsilon>0$ and $t,u\in[0,1]$ let
\begin{gather*}
G_\varepsilon(t,u)=e^{-(t-u)^2/\varepsilon},\qquad m_\varepsilon(t)=\int_0^1G_\varepsilon(t,u)\,du,\\
K_\varepsilon=\frac{G_\varepsilon}{m_\varepsilon(t)},\qquad x_\varepsilon(t)=\int_0^1K_\varepsilon(t,u)x(u)\,du .
\end{gather*}
Since $\int K_\varepsilon(t,u)\,du=1$, the column $(1,x_\varepsilon(t))^T$ equals $\int K_\varepsilon(t,u)(1,x(u))^T\,du$.

\emph{Total positivity.} $G_\varepsilon(t,u)=e^{-t^2/\varepsilon}e^{2tu/\varepsilon}e^{-u^2/\varepsilon}$. The kernel $e^{2tu/\varepsilon}$ is
strictly totally positive (generalized Vandermonde), and positive row and column factors preserve this. Hence
$\det[K_\varepsilon(t_i,u_j)]_{i,j=0}^d>0$ for $t_0<\dots<t_d$ and $u_0<\dots<u_d$; see Karlin, \emph{Total Positivity}, Ch.~3.

\emph{Basic composition formula.} By multilinearity, then antisymmetry in the $u$'s,
\[
\Delta(x_\varepsilon(t_0),\dots,x_\varepsilon(t_d))=\int_{0<u_0<\dots<u_d<1}\det[K_\varepsilon(t_i,u_j)]\;\Delta(x(u_0),\dots,x(u_d))\,du.
\]
The integrand is $\ge0$. Since $x([0,1])$ spans $\R^d$ affinely, some ordered $u^*$ has $\Delta(x(u^*))\ne0$, hence $>0$, and by continuity
$\Delta(x(u))>0$ near $u^*$. Therefore the integral is $>0$, and $x_\varepsilon$ is $d$-order and $C^\infty$.

\emph{Convergence.} $|x_\varepsilon(t)-x(t)|\le\omega(\eta)+2\|x\|_\infty\int_{|u-t|>\eta}K_\varepsilon(t,u)\,du$, where $\omega$ is the modulus of continuity.
Since $m_\varepsilon(t)\ge\frac12\sqrt{\pi\varepsilon}(1-o(1))$ uniformly in $t\in[0,1]$, the last integral is at most
$Ce^{-\eta^2/\varepsilon}/\sqrt\varepsilon\to0$ uniformly. If $x$ is $L$-Lipschitz, then using $\int\partial_tK_\varepsilon\,du=0$,
$|x_\varepsilon'(t)|\le L\int|\partial_tK_\varepsilon(t,u)||u-t|\,du$. Writing
$\partial_tK_\varepsilon=K_\varepsilon\cdot(-2(t-u)/\varepsilon-m_\varepsilon'/m_\varepsilon)$, this is
$\le L\big(2\int K_\varepsilon(u-t)^2/\varepsilon+|m'_\varepsilon/m_\varepsilon|\int K_\varepsilon|u-t|\big)$. Here $|m_\varepsilon'|=|e^{-t^2/\varepsilon}-e^{-(1-t)^2/\varepsilon}|\le1$,
$\int K_\varepsilon(u-t)^2\le C\varepsilon$ and $\int K_\varepsilon|u-t|\le C\sqrt\varepsilon$, so the bound is $O(L)$ uniformly in $\varepsilon\le1$ and $t$.
For the last claim, let $[a,b]\subset[0,1]$. The restrictions $x_\varepsilon|_{[a,b]}$ are $d$-order with the same Lipschitz bound and converge uniformly to
$x|_{[a,b]}$. So (H) on $[a,b]$ follows from Theorem~\ref{thm:ALM} and the two continuity facts below.
(i) $x_n\to x$ uniformly implies $\conv x_n([a,b])\to\conv x([a,b])$ in the Hausdorff metric, and volume is continuous on convex bodies.
(ii) If $x_n\to x$ uniformly with $\operatorname{Lip}(x_n)\le L$, then every iterated integral
$I^{(k)}_n(t)=\int_{a<t_1<\dots<t_k<t}dx_n^{i_1}\cdots dx_n^{i_k}$ converges uniformly to the corresponding $I^{(k)}(t)$ for $x$.
By induction on $k$: for $k=1$ this is clear. If $I^{(k-1)}_n\to I^{(k-1)}$ uniformly, then
$I_n^{(k)}-I^{(k)}=\int(I^{(k-1)}_n-I^{(k-1)})\,dx_n^{i_k}+\int I^{(k-1)}\,d(x_n^{i_k}-x^{i_k})$.
The first term is at most $L\|I^{(k-1)}_n-I^{(k-1)}\|_\infty$. In the second, $I^{(k-1)}$ is Lipschitz, and integrating by parts it equals
$I^{(k-1)}(t)(x_n-x)^{i_k}(t)-\int(x_n-x)^{i_k}\,dI^{(k-1)}$, which tends to $0$ uniformly. Hence $\alpha^{(d)}(x_n|_{[a,b]})\to\alpha^{(d)}(x|_{[a,b]})$. \qed

\section{Proof of Theorem~\ref{thm:planar}}

(c)$\Rightarrow$(b) holds because weak order is preserved under limits.
(b)$\Rightarrow$(c),(a): if $x([0,1])$ spans $\R^2$, apply Proposition~\ref{prop:smooth}. Otherwise $x$ lies on a line $p+\R v$.
A simple closed curve cannot lie on a line, so $x(t)=p+u(t)v$ with $u$ strictly monotone. Put $s=+1$ if $u$ is increasing and $s=-1$ if it is decreasing, and let $J$ be
rotation by $\pi/2$. The paths $x_\varepsilon=p+uv+s\varepsilon u^2Jv$ are $2$-order (they traverse a parabola counterclockwise) and converge to
$x$ in the Lipschitz norm, so $x$ is cyclic. (H) holds trivially, since both sides vanish on a line.
It remains to prove (a)$\Rightarrow$(b).

Both sides of (H) and the sign of $\Delta$ are invariant under increasing absolutely continuous reparametrizations.
Since $x$ is injective on $[0,1)$, the arc-length function is strictly increasing. So we may assume $|x'|=1$ a.e.

\emph{Step 1 (growth at the right end).} Fix $a$. The function $b\mapsto A(a,b)$ is Lipschitz with derivative $\frac12\det(x(b)-x(a),x'(b))$
for a.e.\ $b$. By (H) it coincides with $b\mapsto V(a,b)$, which is nondecreasing. Hence
\begin{equation}\label{eq:i}
\det\big(x(b)-x(a),x'(b)\big)\ge0\quad\text{for every }a\text{ and a.e. }b\in(a,1].
\end{equation}

\emph{Step 2 (growth at the left end).} Fix $b$. Write $A(a,b)=\frac12\int_a^b\det(x,x')\,dt-\frac12\det(x(a),x(b))$.
For a.e.\ $a$, $\partial_aA(a,b)=-\frac12\det(x(a),x'(a))-\frac12\det(x'(a),x(b))=-\frac12\det\big(x'(a),x(b)-x(a)\big)$.
Since $a\mapsto V(a,b)$ is nonincreasing,
\begin{equation}\label{eq:ii}
\det\big(x'(a),x(b)-x(a)\big)\ge0\quad\text{for every }b\text{ and a.e. }a\in[0,b).
\end{equation}

\emph{Step 3 (good base points).} Let $E=\{(a,b):a<b,\ x\text{ differentiable at }a,\ \det(x'(a),x(b)-x(a))<0\}$.
By \eqref{eq:ii} every $b$-section of $E$ is null, so $|E|=0$ by Fubini, and for a.e.\ $a$ the $a$-section is null.
Let $G$ be the set of $a\in[0,1)$ at which $x$ is differentiable, $|x'(a)|=1$, and the $a$-section of $E$ is null.
Then $G$ has full measure. For $a\in G$, continuity of $x$ upgrades \eqref{eq:ii} to all $b>a$: every $x(b)$, $b>a$, lies in the closed half-plane
$H_a=\{w:\det(x'(a),w-x(a))\ge0\}$.

\emph{Step 4 (angle monotonicity).} Fix $a\in G$, $a>0$, and put $y(t)=x(t)-x(a)$. By injectivity $y(t)\ne0$ for $t\in(a,1]$.
Let $\theta$ be a continuous lift of $\arg y$ on $(a,1]$. On each $[a+\eta,1]$, $|y|$ is bounded below, so $\theta$ is Lipschitz there
and $\theta'=\det(y,y')/|y|^2\ge0$ a.e.\ by \eqref{eq:i}. So $\theta$ is nondecreasing.
By Step 3, $e^{i\theta(t)}$ lies in the closed half-circle $\{w\in S^1:\det(x'(a),w)\ge0\}$, whose preimage under $e^{i\cdot}$ is a disjoint
union of closed intervals of length $\pi$. Since $\theta((a,1])$ is connected, it lies in one of them.
Hence for $a<t_1<t_2\le1$ we have $0\le\theta(t_2)-\theta(t_1)\le\pi$ and
\[
\Delta(x(a),x(t_1),x(t_2))=|y(t_1)||y(t_2)|\sin(\theta(t_2)-\theta(t_1))\ge0.
\]

\emph{Step 5 (all triples).} Let $t_0<t_1<t_2$. Choose $a_n\in G\cap(0,t_1)$ with $a_n\to t_0$; this is possible since $G$ has full measure.
By Step 4, $\Delta(x(a_n),x(t_1),x(t_2))\ge0$, and continuity gives $\Delta(x(t_0),x(t_1),x(t_2))\ge0$.
If $x(0)=x(1)$, the triples with $t_0=0$, $t_2=1$ have $\Delta=0$, and the argument above never needs $y(1)\ne0$ for $a=0$. \qed

\begin{remark}[Where injectivity enters, and the role of the source paper]
Steps 1--2 are valid for every path. They are exactly the statement that the hull grows at an endpoint at the rate prescribed by the
signature. Injectivity is used only in Step 4, to exclude returns to the base point. The counterexample of Theorem~\ref{thm:counter}
returns to $x(0)$ and then retraces a straight edge: this sweeps no area and does not enlarge the hull.
The direction (c)$\Rightarrow$(a) is the volume formula for curves with totally positive torsion of de Dios Pont--Ivanisvili--Madrid
\cite{DIM}, Cor.~2.7, extended to $d$-order and cyclic curves in \cite{ALM}. For $d$-order curves the volume formula goes back to Karlin--Studden \cite{KS}; see also
Nudel'man \cite{N} on convex curves and polygonal lines. Proposition~\ref{prop:smooth} complements it by showing that
weak order, a condition on $(d+1)$-point determinants, is all that is needed. The smoothing is a total-positivity argument in the
spirit of Lemma~3.1 of \cite{DIM}.
\end{remark}

\begin{remark}[On Question~\ref{q:inj}]
In $d\ge3$, Step 1 only gives $\langle S^{(d-1)}(x|_{[a,b]})\wedge x'(b),e_1\wedge\dots\wedge e_d\rangle\ge0$ a.e., and similarly at the left end; this is much weaker than in
the plane. The examples of Theorem~\ref{thm:inj} use the fact that a cone needs (H) for its base only on initial intervals. A naive injective lift of the planar example
of Theorem~\ref{thm:counter} (tail $t\mapsto(t-4,0,\varepsilon(t-4))$) turns out to be weakly $3$-ordered, hence cyclic.
\end{remark}

\begin{thebibliography}{9}
\bibitem{ALM} C.~Am\'endola, D.~Lee, C.~Meroni, \emph{Convex hulls of curves: volumes and signatures}, in: Geometric Science of
Information (GSI 2023), LNCS 14071, Springer, 2023, 455--464; arXiv:2301.09405.
\bibitem{DIM} J.~de Dios Pont, P.~Ivanisvili, J.~Madrid, \emph{A new proof of the description of the convex hull of space curves with
totally positive torsion}, Michigan Math.\ J.\ 76 (2026), 1043--1089.
\bibitem{K} S.~Karlin, \emph{Total Positivity}, Vol.~I, Stanford Univ.\ Press, 1968.
\bibitem{KS} S.~Karlin, W.~J.~Studden, \emph{Tchebycheff Systems: With Applications in Analysis and Statistics}, Interscience, 1966.
\bibitem{N} A.~A.~Nudel'man, \emph{Isoperimetric problems for the convex hulls of polygonal lines and curves in multidimensional spaces}, Math.\ USSR Sb.\ 25 (1975) 276--294.
\bibitem{Z} R.~Zielke, \emph{Discontinuous \v Ceby\v sev Systems}, Lecture Notes in Math.\ 707, Springer, 1979.
\end{thebibliography}
\end{document}
